\documentclass[11pt]{article}

\usepackage[T1]{fontenc}
\usepackage[utf8]{inputenc}
\usepackage{lmodern}
\usepackage{amsmath,amssymb,amsfonts,amsthm}
\usepackage[margin=1in]{geometry}
\usepackage{microtype}
\usepackage{xcolor}
\usepackage[colorlinks=true,linkcolor=teal,urlcolor=teal]{hyperref}

\newcommand{\F}{\mathcal F_{\mathrm{orb}}}
\newcommand{\T}{\mathcal T_1}
\newcommand{\Tr}{\operatorname{Tr}}
\newcommand{\dd}{\,\mathrm d}
\newcommand{\B}{\mathrm B}
\newcommand{\Borb}{B_{\mathrm{orb}}}

\newtheorem*{propC}{Proposition C.4 (rephrased)}

\title{A Short Guide to Proposition C.4\\
\large The orbital Gaussian flat-prior upper bound}
\author{}
\date{}

\begin{document}
\maketitle

\begin{abstract}
Proposition C.4 proves that no channel can amplify the displaced thermal
Gaussian model with asymptotic flat-prior fidelity larger than
$\Borb(\gamma,p)$.  Its proof averages a channel over a large phase-space
cube, extracts an exactly covariant trace-nonincreasing limit using only
compact observables, and evaluates a specially chosen weighted moment.  The
compact witness avoids any need to pass fidelity itself through the weak
limit.  This note explains the argument and how Lemma C.1, Lemma C.2, and
Proposition C.3 fit together.
\end{abstract}

\section{The Gaussian problem}

The orbital LAN model consists of displaced product thermal states on a
finite-mode Fock space $\F$:
\[
  \Phi_z=D(z)\tau D(z)^*,
  \qquad
  \Psi_z=D(\sqrt\gamma z)\tau D(\sqrt\gamma z)^*.
\]
The notation suppresses that the target has the same centered thermal state
$\tau$ as the ideal product model, whereas a physical amplifier necessarily
produces the noisier centered state $\tau^{(\gamma)}$.  The product
quantum-limited amplifier sends
\[
  \Phi_z\longmapsto
  D(\sqrt\gamma z)\tau^{(\gamma)}D(\sqrt\gamma z)^*.
\]
Its fidelity with $\Psi_z$ is independent of $z$ and equals
\[
  \B(\tau^{(\gamma)},\tau)=\Borb(\gamma,p).
\]

For the cube
\[
  \mathsf Z_L
  =\{z:|\Re z_\ell|,|\Im z_\ell|\leq L\},
\]
Proposition C.4 supplies the converse inequality
\begin{equation}\label{eq:goal}
  \limsup_{L\to\infty}\sup_\Lambda
  \frac1{|\mathsf Z_L|}\int_{\mathsf Z_L}
  \B(\Lambda(\Phi_z),\Psi_z)\dd z
  \leq\Borb(\gamma,p).
\end{equation}

\section{The witness and its three properties}

Write the two centered thermal states in their joint number basis:
\[
  \tau=\sum_{\mathbf k}t_{\mathbf k}|\mathbf k\rangle\!\langle\mathbf k|,
  \qquad
  \tau^{(\gamma)}
  =\sum_{\mathbf k}t_{\mathbf k}^{(\gamma)}
    |\mathbf k\rangle\!\langle\mathbf k|.
\]
Define
\begin{equation}\label{eq:witness}
  W_\gamma
  =\sum_{\mathbf k}
  \sqrt{\frac{t_{\mathbf k}}{t_{\mathbf k}^{(\gamma)}}}
  |\mathbf k\rangle\!\langle\mathbf k|.
\end{equation}
This witness is chosen so that three otherwise separate parts of the proof
match exactly.

\begin{enumerate}
  \item It is bounded, positive, injective, and compact.  Its eigenvalues are
  products of decreasing geometric sequences and tend to zero at high total
  photon number.
  \item It is an admissible product of decreasing photon-number weights for
  the least-noise moment in Proposition C.3.
  \item It satisfies
  \begin{equation}\label{eq:witness-id}
    \Tr(\tau^{(\gamma)}W_\gamma)
    =\Tr(\tau W_\gamma^{-1})
    =\sum_{\mathbf k}\sqrt{t_{\mathbf k}t_{\mathbf k}^{(\gamma)}}
    =\Borb(\gamma,p).
  \end{equation}
\end{enumerate}

Lemma C.1, the weighted fidelity inequality, therefore gives for every
positive trace-class $R$
\begin{equation}\label{eq:weighted-B}
  \B(R,\tau)^2
  \leq \Tr(RW_\gamma)\,\Borb(\gamma,p).
\end{equation}
The nonlinear fidelity has been reduced to one linear compact-observable
moment.

\section{Statement of Proposition C.4}

\begin{propC}
For the orbital Gaussian input and target families,
\[
  \limsup_{L\to\infty}\sup_\Lambda
  \int_{\mathsf Z_L}\B(\Lambda(\Phi_z),\Psi_z)
  \frac{\dd z}{|\mathsf Z_L|}
  \leq\Borb(\gamma,p),
\]
where the supremum is over all channels on the orbital Fock space.
\end{propC}

\section{Proof walkthrough}

\subsection*{Step 1: move every displacement back to the origin}

For a channel $\Lambda$, define its finite F\o lner average
\begin{equation}\label{eq:average}
  \overline\Lambda_L
  =\frac1{|\mathsf Z_L|}\int_{\mathsf Z_L}
  \operatorname{Ad}_{D(-\sqrt\gamma z)}\circ\Lambda\circ
  \operatorname{Ad}_{D(z)}\,\dd z.
\end{equation}
It is again a channel.  Joint concavity and unitary invariance of root
fidelity imply
\begin{equation}\label{eq:concavity}
  \frac1{|\mathsf Z_L|}\int_{\mathsf Z_L}
  \B(\Lambda(\Phi_z),\Psi_z)\dd z
  \leq\B(\overline\Lambda_L(\tau),\tau).
\end{equation}
Thus all parameter values are collected into the performance of one averaged
channel at the centered input.

\subsection*{Step 2: the averaged maps become covariant}

Choose $L_j\to\infty$ and channels whose payoffs approach the limsup in
\eqref{eq:goal}; write $\Gamma_j=\overline\Lambda_{L_j}$.  Shifting the
integration variable in \eqref{eq:average} gives, for every fixed $w$,
\begin{align}
 &\|\Gamma_j(D(w)XD(w)^*)
 -D(\sqrt\gamma w)\Gamma_j(X)D(\sqrt\gamma w)^*\|_1\notag\\
 &\qquad\leq
 \frac{|(\mathsf Z_{L_j}+w)\mathbin\triangle\mathsf Z_{L_j}|}
      {|\mathsf Z_{L_j}|}\,\|X\|_1.
 \label{eq:boundary}
\end{align}
Only the boundary portion on which the shifted and unshifted cubes differ
contributes.  Euclidean cubes are a F\o lner family, so the ratio on the
right tends to zero.

Lemma C.2 now yields a subsequence and a completely positive
trace-nonincreasing map $\Gamma$ such that
\begin{equation}\label{eq:compact-limit}
  \Tr[\Gamma_j(X)K]\longrightarrow\Tr[\Gamma(X)K]
  \qquad(K\text{ compact}).
\end{equation}
Equation \eqref{eq:boundary} also makes the limit exactly displacement
covariant with gain $\sqrt\gamma$.

The limit may lose trace because oscillator mass can escape to high energy,
but this causes no problem: Proposition C.3 was formulated for
trace-nonincreasing maps.

\subsection*{Step 3: pass the compact witness moment to the limit}

Apply \eqref{eq:weighted-B} to $R=\Gamma_j(\tau)$:
\[
  \B(\Gamma_j(\tau),\tau)^2
  \leq
  \Tr[\Gamma_j(\tau)W_\gamma]\,\Borb(\gamma,p).
\]
Since $W_\gamma$ is compact, \eqref{eq:compact-limit} gives
\begin{equation}\label{eq:moment-limit}
  \Tr[\Gamma_j(\tau)W_\gamma]
  \longrightarrow
  \Tr[\Gamma(\tau)W_\gamma].
\end{equation}
This is the key reason for introducing the weighted witness before taking a
limit.  One never needs an upper-semicontinuity theorem for fidelity in the
chosen topology.

\subsection*{Step 4: use the least-noise moment}

The limiting map has exactly the hypotheses of Proposition C.3.  Therefore
\begin{align*}
  \Tr[\Gamma(\tau)W_\gamma]
  &\leq
  \Tr[\Gamma(\tau)]
  \Tr[\tau^{(\gamma)}W_\gamma]\\
  &\leq\Borb(\gamma,p),
\end{align*}
where the last line uses $\Tr\Gamma(\tau)\leq1$ and
\eqref{eq:witness-id}.  Combining this with
\eqref{eq:weighted-B}--\eqref{eq:moment-limit} gives
\[
  \limsup_j\B(\Gamma_j(\tau),\tau)^2
  \leq\Borb(\gamma,p)^2.
\]
Root fidelity is nonnegative, so taking square roots and returning to
\eqref{eq:concavity} proves \eqref{eq:goal}.

\section{Why the argument is sharp}

The witness is not an arbitrary estimate.  Its eigenvalue at $\mathbf k$ is
the square-root likelihood ratio between the ideal thermal law and the
quantum-limited output law.  Consequently, both factors in the weighted
Cauchy--Schwarz bound equal their classical Hellinger overlap.  The
least-noise proposition then says that no admissible covariant limit can have
a larger expectation of this witness.

The product quantum-limited amplifier attains $\Borb(\gamma,p)$ at every
$z$, so Proposition C.4 supplies the matching upper bound used in the
Gaussian flat-prior theorem.

\section{Logical dependencies}

The proof has the compact chain
\[
  \begin{array}{c}
  \text{F\o lner averaging and concavity}\\
  \Downarrow\\
  \text{Lemma C.2: compact-observable covariant limit}\\
  \Downarrow\\
  \text{Lemma C.1: fidelity }\leq\text{ witness moment}\\
  \Downarrow\\
  \text{Proposition C.3: least-noise moment}\\
  \Downarrow\\
  \Borb(\gamma,p).
  \end{array}
\]
The proof is entirely orbital: there is no classical spectral coordinate.
Proposition C.6 repeats this architecture with an additional classical
witness and convolution argument.

\end{document}
